% !TeX root = ../../Book4.tex
% 纲要标题：### 5.2 带边界单连通区域的映照
% 纲要位置：几何学计划纲要.md，第 3641 行。
% 正文状态：初稿；尚未正式审核。

% 纲要内容开始
% 
% * 拓扑圆盘和边界正则化
% * 由 Green 函数构造全纯映射
% * 边界模长与映射次数
% * 单叶性和圆盘映照
% 
% 纲要内容结束

\section{带边界单连通区域的映照}
\label{g4:sec:0502}

在一个没有预先给定平面坐标的曲面区域上，Green 函数仍能测量点到边界的势差。它在极点附近像 \(-\log|z|\)，而边界值为零。如果把它补成局部全纯函数，再取负指数，极点就变成一个简单零点，边界则变成模长为 \(1\) 的限制。下面证明这个构造给出整个单位圆盘的单叶坐标。单连通性用于消除指数的多值性；边界条件用于保证原像不逃向边界。

\subsection{拓扑圆盘和边界正则化}
\label{g4:subsec:0502:01}

设 \(X\) 为黎曼曲面。称紧子集 \(D\subset X\) 为嵌入的闭拓扑圆盘，是指存在从闭单位圆盘到 \(D\) 的同胚，且它把开圆盘送到 \(D\) 在 \(X\) 中的内部。以后写 \(U=\operatorname{int}D\)。若 \(\partial D\) 是光滑嵌入圆周，则 \(D\) 同时是紧光滑带边界曲面，\(U\Subset X\)。这里的拓扑条件只描述区域的形状，并没有提供全纯坐标；把它变成解析圆盘正是本节的目标。

\ProseParagraphBreak

由有限条坐标弧围成的盘常带有角点。势论中的光滑边界条件可以通过局部修改实现，并且修改能够避开事先选定的内部紧集。下述引理只处理已经得到的拓扑盘；下一节另行证明一般非紧单连通曲面如何被这样的盘穷竭。

\begin{lemma}
	\label{g4:lem:0502:round-corners}
	设 \(D\) 是黎曼曲面 \(X\) 中的嵌入闭拓扑圆盘，其边界为分片光滑正则曲线，只有有限个角点，且每个角点处的两条单侧切线横截。给定紧集 \(K\subset\operatorname{int}D\) 及包含 \(D\) 的开集 \(W\)，存在光滑边界的嵌入闭拓扑圆盘 \(D'\subset W\)，满足 \(K\subset\operatorname{int}D'\)。修改可以限制在各角点任意预先选定的充分小邻域内。
\end{lemma}
\begin{proof}
	取彼此不交、避开 \(K\) 且闭包包含于 \(W\) 的角点坐标邻域。只需说明一个角点附近的构造，因为随后可以在这些不交邻域中分别实施。

	将角点坐标平面中的坐标轴作实线性变换后，两条单侧弧都能以横坐标为参数，并分别位于横坐标的两侧。因此在一个小矩形中，边界写成 \(y=b(x)\)：函数 \(b\) 连续、分段光滑，在 \(x=0\) 处只有一阶导数的跳跃；区域在这张图像的一侧。这样的坐标方向存在，因为两条离开角点的切向射线不同且不共线，可以选择一个线性泛函，使它在两条射线上取相反符号。缩小矩形后，两弧在这个方向上的导数都不为零，便得到所需图像。

	把 \(b\) 在稍大的区间上延拓为连续函数，取光滑非负、积分为 \(1\) 的紧支核，作卷积 \(b*\varphi_\varepsilon\)。卷积在紧区间上一致逼近 \(b\)。选光滑截断函数 \(\chi(x)\)，在 \(0\) 附近为 \(1\)，在上述小区间端部为 \(0\)，令
	\[
		b_\varepsilon(x)=\chi(x)(b*\varphi_\varepsilon)(x)
		+(1-\chi(x))b(x).
	\]
	这是光滑函数：在唯一可能不光滑的点 \(0\) 附近只剩卷积项；在截断过渡区，原函数本来光滑。它与 \(b\) 在端部相等，并一致趋于 \(b\)。

	还需保证换图像不会改变盘的拓扑。记 \(\delta(x)=b_\varepsilon(x)-b(x)\)，取一个光滑函数 \(\beta\)，在 \(0\) 附近为 \(1\)，支集放在很短的纵向区间内。在矩形中定义
	\[
		T(x,y)=\bigl(x,y+\beta(y-b(x))\delta(x)\bigr).
	\]
	取 \(\varepsilon\) 充分小，使 \(\|\beta'\|_\infty\|\delta\|_\infty<1/2\)。对每个固定 \(x\)，第二分量关于 \(y\) 严格递增，并在区间端部等于 \(y\)，因而是该纵向区间到自身的同胚。它的斜率至少为 \(1/2\)，所以逆映射连续依赖 \(x,y\)。于是 \(T\) 是矩形的同胚，在边界附近恒等，且把旧图像送到新图像。将它在矩形外延为恒等映射，便得到 \(X\) 的一个同胚，其像盘的边界在该角点处已经光滑。

	对所有角点作上述修改。所得环境同胚在 \(K\) 附近恒等，修改支集包含在 \(W\) 内，故 \(D'=T(D)\) 仍为闭拓扑盘，包含 \(K\) 于内部且包含于 \(W\)。边界在修改区为光滑图像，在接缝附近与旧光滑弧完全相同，所以全边界光滑。
\end{proof}

这个构造改变的是区域边界，曲面已有的复结构始终保留。因而可以在正则化后的区域直接使用\cref{g4:thm:0501:dirichlet,g4:thm:0501:green}。只要极点与需要控制的紧集留在内部，后面的映照构造就适用。

\subsection{由 Green 函数构造全纯映射}
\label{g4:subsec:0502:02}

现在固定光滑边界闭拓扑盘 \(D\subset X\)，令 \(U=\operatorname{int}D\)，并选 \(p\in U\)。取 \(p\) 附近的全纯坐标 \(z\)，满足 \(z(p)=0\)。上一节构造了 Green 函数 \(G=G_U(\cdot,p)\)，它在 \(U\setminus\{p\}\) 调和且为正，连续延至 \(\partial U\) 取零值，并在 \(p\) 附近满足
\begin{equation}
	\label{g4:eq:0502:green-local}
	G=-\log|z|+h,
\end{equation}
其中 \(h\) 在整个小坐标盘内调和。

\ProseParagraphBreak

依照上一节的约定，\(*\mathrm{d}x=\mathrm{d}y\)、\(*\mathrm{d}y=-\mathrm{d}x\)。因为 \(G\) 调和，\(*\mathrm{d}G\) 在去点区域上是闭一形式。它的局部原函数 \(v\) 满足 \(\mathrm{d}v=*\mathrm{d}G\)，所以 \(G+\mathrm{i}v\) 局部全纯。绕极点一周时，\(v\) 不会回到原值；需要先算出这个变化，再判断指数能否成为单值函数。

\begin{lemma}
	\label{g4:lem:0502:green-periods}
	设 \(U\) 是光滑边界闭拓扑盘的内部，\(p\in U\)，\(G=G_U(\cdot,p)\)。对每条分片光滑闭路 \(\gamma\subset U\setminus\{p\}\)，有
	\[
		\int_\gamma *\mathrm{d}G\in2\pi\mathbb Z.
	\]
	若 \(\gamma\) 是以 \(p\) 为中心的全纯坐标中的小圆周，取逆时针方向，则该积分为 \(-2\pi\)。
\end{lemma}
\begin{proof}
	先计算小圆周。由式\eqref{g4:eq:0502:green-local}，
	\[
		*\mathrm{d}G=-\mathrm{d}\arg z+*\mathrm{d}h.
	\]
	\(h\) 在整个坐标盘内调和，故在该盘内有单值调和共轭，后一项是恰当形式。沿逆时针小圆周积分，结果为 \(-2\pi\)。

	说明这个局部计算为何控制全部闭路。以 \(p\) 的一个小坐标圆盘 \(B\) 覆盖缺失点，便有开覆盖
	\[
		U=(U\setminus\{p\})\cup B,
		(U\setminus\{p\})\cap B=B\setminus\{p\}.
	\]
	这三个开集均道路连通。\(B\setminus\{p\}\) 径向形变收缩到小圆周，其基本群由该圆周生成，这是第一卷第5.1节定理5.1.13的圆周绕数计算。第一卷第5.2节定理5.2.4的 van Kampen 定理\footnote{Egbert Rudolf van Kampen（1908--1942），生于比利时的数学家，在荷兰求学、美国任教，研究拓扑与拓扑群。}说明：因为 \(B\) 和 \(U\) 都单连通，\(\pi_1(U\setminus\{p\})\) 由这条小圆周的共轭及其逆元正规生成。

	由\cref{g4:lem:0501:closed-form-periods}，闭一形式的积分只依赖定端点连续同伦类，并对道路拼接可加；因此它给出从基本群到实数加法群的同态。沿共轭道路 \(\alpha*\gamma*\bar\alpha\) 的积分等于沿 \(\gamma\) 的积分，连接弧的贡献相消。所有闭路的周期于是都是小圆周周期 \(-2\pi\) 的整数倍。
\end{proof}

\begin{proposition}
	\label{g4:prop:0502:green-map}
	设 \(D\) 是黎曼曲面中的光滑边界闭拓扑盘，\(U=\operatorname{int}D\)，\(p\in U\)。存在全纯函数 \(F:U\to\mathbb D\)，满足
	\[
		F(p)=0,\qquad |F(q)|=e^{-G_U(q,p)}\quad(q\ne p).
	\]
	它只在 \(p\) 有零点，且该零点为简单零点；这样的函数彼此只差一个模为 \(1\) 的常数因子。若在中心坐标 \(z\) 中写 \(G=-\log|z|+h\)，则
	\[
		\left|\frac{\mathrm{d}F}{\mathrm{d}z}(p)\right|=e^{-h(p)}.
	\]
\end{proposition}
\begin{proof}
	在 \(U\setminus\{p\}\) 中选基点 \(q_0\)，先给局部共轭指定一个实常数。对从 \(q_0\) 到 \(q\) 的道路积分 \(*\mathrm{d}G\)，得到一个可能依赖道路的数 \(v(q)\)。由\cref{g4:lem:0502:green-periods}，改变道路使这个数增加 \(2\pi\) 的整数倍。因此
	\begin{equation}
		\label{g4:eq:0502:exponential-map}
		F(q)=\exp\bigl(-G(q)-\mathrm{i}v(q)\bigr)
	\end{equation}
	与道路无关。在任一不含 \(p\) 的小坐标盘上，\(G+\mathrm{i}v\) 是全纯函数，因而 \(F\) 全纯且不为零。\(G>0\) 给出 \(|F|<1\)。

	在 \(p\) 附近，给 \(h\) 选单值调和共轭 \(h^*\)，使 \(H=h+\mathrm{i}h^*\) 全纯。在一个割开的小盘中，\(v=-\arg z+h^*+c\)，其中 \(c\) 为实常数。式\eqref{g4:eq:0502:exponential-map}便成为
	\[
		F=e^{-\mathrm{i}c}\,z\,e^{-H(z)}.
	\]
	右边在完整小盘上全纯，因此补上 \(F(p)=0\) 后得到 \(U\) 上的全纯函数；该式还说明零点恰为一阶，且导数模为 \(e^{-h(p)}\)。

	若 \(F_1,F_2\) 都有上述模长，二者在 \(p\) 的简单零点相消，故 \(F_1/F_2\) 在 \(U\) 上处处全纯，模恒为 \(1\)。最大模原理使这个比值为常数。反之，乘任一模为 \(1\) 的常数显然保留全部要求。
\end{proof}

\subsection{边界模长与映射次数}
\label{g4:subsec:0502:03}

Green 边界值给出 \(|F|\to1\)，这已经足以控制所有位于圆盘内部的原像。确实，\(|F|\) 在 \(\overline U\) 上连续延拓，边界值为 \(1\)。对任意 \(0<r<1\)，集合
\[
	\{q\in U:|F(q)|\le r\}
\]
是 \(\overline U\) 中避开边界的闭集，因此紧且包含在 \(U\) 内。目标中的每个紧集都包含在某个 \(\{|w|\le r\}\) 中，其原像便紧。也就是说，\(F:U\to\mathbb D\) 是固有映射。这里控制的是模长；证明不要求先把 \(F\) 的相位逐点延拓到边界。

\ProseParagraphBreak

第二章的次数定理以紧源曲面为假设。当前的 \(U\) 是开曲面，所需的紧性由固有性提供。下面将对应的计数论证写出，使这个变化的作用明确可见。

\begin{proposition}
	\label{g4:prop:0502:proper-degree}
	设 \(S,T\) 是连通黎曼曲面，\(f:S\to T\) 是固有的非恒定全纯映射。则 \(f\) 满射，每个纤维有限，且存在整数 \(d\ge1\)，使对每个 \(y\in T\) 都有
	\[
		\sum_{x\in f^{-1}(y)}e_x(f)=d.
	\]
	这个整数称为 \(f\) 的次数。若某个纤维只有一个点，且该点的分歧指数为 \(1\)，则 \(f\) 为双全纯映射。
\end{proposition}
\begin{proof}
	由\cref{g4:thm:0202:local-model}，\(f\) 的局部形式为 \(z\mapsto z^e\)，所以它开，纤维中的各点也孤立。固有性使纤维紧；紧离散空间有限，所以每个纤维有限。

	像还在 \(T\) 中闭。若像中的点列 \(y_n\) 趋于 \(y\)，可以把尾部与 \(y\) 放在同一个紧坐标闭盘内。任选 \(x_n\in f^{-1}(y_n)\)，固有性使这些原像位于同一个紧集。取收敛子列 \(x_{n_k}\to x\)，连续性给出 \(f(x)=y\)。黎曼曲面有可数局部基，这个点列判据足以证明闭性。非空的像既开又闭，而 \(T\) 连通，故 \(f\) 满射。

	固定 \(y\)，记纤维为 \(x_1,\ldots,x_m\)。在各点取闭包互不相交的小坐标盘 \(B_i\)，使 \(f\) 的局部幂模型成立，并使 \(x_i\) 是 \(\overline B_i\) 中 \(y\) 的唯一原像。充分接近 \(y\) 的目标点，其全部原像都在这些盘内：否则可取 \(y_n\to y\) 和 \(x_n\notin\bigcup_iB_i\) 满足 \(f(x_n)=y_n\)，固有性再次给出收敛子列，其极限既在这个开集的补集中，又属于 \(f^{-1}(y)\)，矛盾。

	目标坐标中取充分小的 \(y'\) 邻域后，局部幂模型说明 \(B_i\) 内全部原像的重数之和恒为 \(e_{x_i}(f)\)；对 \(y'\ne y\)，这些原像实际上互异且各重数为 \(1\)。因此总重数关于 \(y\) 局部常值。\(T\) 连通，得到固定的正整数 \(d\)。若某个纤维总重数为 \(1\)，则每个纤维都恰有一个点，而且处处分歧指数为 \(1\)。于是 \(f\) 双射，局部全纯逆又按唯一性拼成全局全纯逆。
\end{proof}

对 Green 函数构造的 \(F\)，零纤维恰为 \(\{p\}\)，且刚已证明 \(e_p(F)=1\)。所以它的次数为 \(1\)。这一步同时得到满射性和单射性：边界条件保证每个值都能取到，而一个简单零点把所有纤维的总重数固定为 \(1\)。

\subsection{单叶性和圆盘映照}
\label{g4:subsec:0502:04}

把上述构造与计数合在一起，就得到适合在穷竭中反复使用的归一化映照。归一化必须说明在哪张源坐标中测量导数，因为曲面上两个不同的坐标会给出不同的导数数值。

\begin{theorem}
	\label{g4:thm:0502:bordered-disk}
	设 \(X\) 为黎曼曲面，\(D\subset X\) 是光滑边界的嵌入闭拓扑圆盘，\(U=\operatorname{int}D\)。给定 \(p\in U\) 及 \(p\) 附近的中心全纯坐标 \(z\)，存在唯一双全纯映射
	\[
		f_U:U\longrightarrow\mathbb D,
		\qquad f_U(p)=0,\qquad
		\frac{\mathrm{d}f_U}{\mathrm{d}z}(p)>0.
	\]
	最后一个不等式指正实数。它满足
	\[
		G_U(q,p)=-\log|f_U(q)|\quad(q\ne p).
	\]
	若 \(G_U=-\log|z|+h\) 于 \(p\) 附近，则 \(f_U'(p)=e^{-h(p)}\)。
\end{theorem}
\begin{proof}
	由\cref{g4:prop:0502:green-map}构造 \(F\)。边界模长使它固有，\cref{g4:prop:0502:proper-degree}与简单零点使它双全纯。乘以唯一的模为 \(1\) 的常数，使 \(F\) 在坐标 \(z\) 中的非零导数成为正实数，便得所求 \(f_U\)，Green 公式与导数公式也随构造成立。

	若 \(f_1,f_2\) 都满足归一化条件，则 \(f_2\circ f_1^{-1}\) 是固定 \(0\) 的圆盘自同构。由\cref{g4:thm:0103:disk-aut}，它为旋转 \(w\mapsto e^{\mathrm{i}\theta}w\)。在基点求导，有 \(e^{\mathrm{i}\theta}=f_2'(p)/f_1'(p)>0\)，从而 \(e^{\mathrm{i}\theta}=1\)，两映射相同。
\end{proof}

\begin{corollary}
	\label{g4:cor:0502:domain-monotonicity}
	设 \(U\subset V\) 是同一黎曼曲面中两个光滑边界闭拓扑盘的内部，\(p\in U\)。在同一个中心坐标中取上述归一化映照 \(f_U,f_V\)。则
	\[
		0<f_V'(p)\le f_U'(p),\qquad
		G_V(q,p)\ge G_U(q,p)\quad(q\in U\setminus\{p\}).
	\]
	导数等号成立当且仅当 \(U=V\)。
\end{corollary}
\begin{proof}
	复合 \(H=f_V\circ f_U^{-1}:\mathbb D\to\mathbb D\) 全纯且固定 \(0\)。Schwarz 引理给出 \(|H(w)|\le|w|\) 及 \(|H'(0)|\le1\)。代入 \(w=f_U(q)\)，再用链式法则，得到
	\[
		|f_V(q)|\le|f_U(q)|,\qquad
		H'(0)=\frac{f_V'(p)}{f_U'(p)}\in(0,1].
	\]
	取负对数便得到 Green 函数的不等式。若导数比为 \(1\)，Schwarz 引理的等号条件使 \(H\) 为恒等映射，故 \(f_V(U)=\mathbb D=f_V(V)\)。\(f_V\) 单射，于是 \(U=V\)。反向由归一化映照的唯一性成立。
\end{proof}

\begin{example}
	\label{g4:exm:0502:disk-green}
	对平面圆盘 \(U=\{|z|<R\}\)，极点 \(p\in U\)，以 \(z-p\) 为中心坐标，归一化映照和 Green 函数分别为
	\[
		f_U(z)=\frac{R(z-p)}{R^2-\overline pz},\qquad
		G_U(z,p)=\log\left|\frac{R^2-\overline pz}{R(z-p)}\right|.
	\]
	分母在 \(U\) 内不为零，且
	\[
		f_U'(p)=\frac{R}{R^2-|p|^2}>0.
	\]
	当 \(p=0\) 时，\(G_U=\log(R/|z|)\)，构造所得坐标正是 \(z/R\)。这些公式依赖区域的共形结构；为定义 Laplace 算子而选的辅助共形度量不出现在答案中。
\end{example}

现在每个光滑相对紧拓扑盘都有唯一的归一化全纯坐标，嵌套区域上的基点导数也有确定的单调方向。要将这个局部于区域的结论推广到整张非紧曲面，还需构造合适的盘穷竭，并证明归一化映照的极限不会退化；下一节分别处理这两个问题。
% 纲要条目排为编号小节标题，并进入目录与书签；说明段仍保留于纲要注释中。